%! Piecewise-Defined Functions — Unit 2, Lesson 2.6 — Homework (Answer Key)
\documentclass[10pt]{article}
\usepackage{atda-article}
\usepackage{atda-key}   % pulls in -boxes; do NOT also load -boxes
\newcommand{\TallMath}[1]{$\displaystyle #1\rule[-1.4em]{0pt}{3.2em}$}
\setlength{\workrowsep}{8pt}

\begin{document}

\pageheader{Unit 2, Lesson 2.6}{Homework --- Answer Key}
% NAMESTRIP: no \namedateperiod here — mirrors the blank. Teacher-only prose goes
% in the lesson plan as \begin{teachernote}[Homework], never in a key.

\vspace{0.15in}

\boxguard[24]
\begin{notesbox}{Practice --- Which Piece Owns This Input?}
\small
Before you use a rule, say which piece owns the input. At a boundary, let the dots decide.

\begin{enumerate}[leftmargin=1.6em, itemsep=6pt, topsep=4pt]

\item \textbf{Reading a piecewise graph.} Use the graph of $f$ below. Fill both rows --- for the
second row write \emph{1st} or \emph{2nd} to say which piece owns that input.

\begin{center}
\begin{tikzpicture}
\begin{axis}[
    width=9.8cm, height=4.4cm,
    xlabel={$x$}, ylabel={$y$},
    xmin=-5, xmax=5, ymin=-3, ymax=6,
    xtick={-5,-4,-3,-2,-1,0,1,2,3,4,5}, ytick={-2,-1,0,1,2,3,4,5},
    grid=both, grid style={linegray!45},
    axis lines=left,
    tick label style={font=\tiny}, label style={font=\scriptsize}]
  \addplot[royal, very thick] coordinates {(-4,-2) (-1,1)};
  \addplot[royal, very thick] coordinates {(-1,5) (4,0)};
  \addplot[royal, only marks, mark=*, mark size=1.9pt] coordinates {(-4,-2) (-1,1) (4,0)};
  \addplot[royal, only marks, mark=*, mark size=1.9pt, mark options={fill=white}]
    coordinates {(-1,5)};
\end{axis}
\end{tikzpicture}
\end{center}

\vspace{-2pt}
\renewcommand{\arraystretch}{1.3}
\begin{tabular}{@{}|l|c|c|c|c|c|@{}}
\hline
$x$ & $-3$ & $-1$ & $0$ & $2$ & $4$ \\ \hline
$f(x)$ & \ans{$-1$} & \ans{$1$} & \ans{$4$} & \ans{$2$} & \ans{$0$} \\ \hline
Which piece? & \ans{1st} & \ans{1st} & \ans{2nd} & \ans{2nd} & \ans{2nd} \\
\hline
\end{tabular}

\item \textbf{The other direction --- from the rules, with no graph.} Let
\[
  g(x) = \begin{cases}
    x+4 & \text{if } -4 \le x \le 0,\\[2pt]
    4 & \text{if } 0 < x \le 3,\\[2pt]
    4-2(x-3) & \text{if } 3 < x \le 6.
  \end{cases}
\]
Fill both rows. Name the piece by its rule number (1st, 2nd or 3rd).
\smallskip

\renewcommand{\arraystretch}{1.3}
\begin{tabular}{@{}|l|c|c|c|c|c|@{}}
\hline
$x$ & $-2$ & $0$ & $2$ & $5$ & $6$ \\ \hline
Which piece? & \ans{1st} & \ans{1st} & \ans{2nd} & \ans{3rd} & \ans{3rd} \\ \hline
$g(x)$ & \ans{$2$} & \ans{$4$} & \ans{$4$} & \ans{$0$} & \ans{$-2$} \\
\hline
\end{tabular}

\item \textbf{Continuous or not?} Each line describes what happens at one boundary point. Write
\emph{continuous} or \emph{discontinuous}, and give the reason in a few words.
\smallskip

\renewcommand{\arraystretch}{1.35}
\begin{tabularx}{\linewidth}{@{}X p{3.4cm}@{}}
\rowcolor{mist}
\textbf{At the boundary\dots} & \textbf{The graph is\dots} \\
\hline
both rules give the value $12$ there & \ans{continuous --- they meet} \\
the left rule gives $5$ and the right rule gives $9$ & \ans{discontinuous --- a jump of $4$} \\
the graph turns a sharp corner, and both pieces reach the same point & \ans{continuous --- a corner
    is not a break} \\
there is a filled circle at $(3,2)$ and an open circle at $(3,7)$ & \ans{discontinuous --- a jump of
    $5$} \\
\end{tabularx}

% Unguarded, item 4's stem and its whole graph sat at the foot of p1 with the
% question and all four write-line slots opening p2 -- a student answering on one
% page while the figure they must read is on the previous one. Same defect class
% as 2.4's item 4 and 2.5's item 8. \boxguard is inert inside a breakable
% tcolorbox, so \tcbbreak is the only lever, and here it costs nothing: still 4
% pages, and it fills p4, which was otherwise ~45% empty. Mirrors homework.
\tcbbreak
\item \textbf{A bike-share priced in steps.} A city bike costs \$2 for a ride of up to and including
$30$ minutes, \$5 for over $30$ up to and including $60$, and \$9 for over $60$ up to and including
$120$.

\begin{center}
\begin{tikzpicture}
\begin{axis}[
    width=9.8cm, height=3.8cm,
    xlabel={$t$ (minutes of the ride)}, ylabel={$R$ (cost, \$)},
    xmin=0, xmax=126, ymin=0, ymax=12,
    xtick={0,15,30,45,60,75,90,105,120}, ytick={0,2,4,6,8,10},
    minor y tick num=1,
    grid=both, grid style={linegray!45}, minor grid style={linegray!30},
    axis lines=left,
    tick label style={font=\tiny}, label style={font=\scriptsize}]
  \addplot[royal, very thick] coordinates {(0,2) (30,2)};
  \addplot[royal, very thick] coordinates {(30,5) (60,5)};
  \addplot[royal, very thick] coordinates {(60,9) (120,9)};
  \addplot[royal, only marks, mark=*, mark size=1.9pt] coordinates {(30,2) (60,5) (120,9)};
  \addplot[royal, only marks, mark=*, mark size=1.9pt, mark options={fill=white}]
    coordinates {(0,2) (30,5) (60,9)};
\end{axis}
\end{tikzpicture}
\end{center}

Give $R(30)$ and $R(60)$, naming the circle that settles each. Then give the range of $R$, and name
every input where you would have to lift your pencil.

\ansline{$R(30)=\$2$ --- the \emph{filled} circle at $(30,2)$; the circle at $(30,5)$ is open. A ride of}\\
\ansline{exactly $30$ minutes still costs \$2. $R(60)=\$5$ --- the filled circle at $(60,5)$, same reason.}\\
\ansline{Range: the list $\{2, 5, 9\}$ --- only three prices are ever charged, so no interval describes it.}\\
\ansline{Lift the pencil at $t=30$ and at $t=60$: two jumps, of \$3 and of \$4.}\\

\item \textbf{The full read.} Go back to the graph of $f$ in item 1 and answer all of it.
\begin{enumerate}[label=(\alph*), leftmargin=1.6em, itemsep=3pt, topsep=3pt]
  \item Domain and range. (One of them needs a parenthesis --- look at the open circle.)
  \item The zeros and the $y$-intercept.
  \item The interval on which $f$ is increasing and the interval on which it is decreasing.
  \item The absolute maximum and the absolute minimum, each with a value and a location. One of the
        two does not exist --- say which, and why.
  \item Is $f$ continuous? If not, name the input where it breaks and give $f$ there.
\end{enumerate}

\ansline{(a) Domain $[-4,4]$; range $[-2,5)$ --- a parenthesis at $5$, because that circle is open.}\\
\ansline{(b) Two zeros, $x=-2$ and $x=4$; $y$-intercept $(0,4)$ --- $x=0$ is owned by the second piece.}\\
\ansline{(c) Increasing on $[-4,-1]$; decreasing on $(-1,4]$.}\\
\ansline{(d) No absolute maximum: just right of $x=-1$ the outputs get as close to $5$ as you like and}\\
\ansline{never reach it. Absolute minimum $-2$ at $x=-4$. \quad (e) Not continuous --- it jumps at $x=-1$, and $f(-1)=1$.}\\

\end{enumerate}
\end{notesbox}

\vspace{0.10in}

% p2 ends ~55% full and this box opens p3. Lowering the guard [30] -> [22] was
% tested and changed nothing at all: the box's first unbreakable chunk (scenario
% paragraph + the 4.4cm figure, ~2.9in) is larger than p2's free space, so the
% break is the natural one, not a guard problem. Restored per 1.5's
% test-then-restore rule. Mirrors homework.
\boxguard[30]
\begin{scenariobox}[In Context --- A Tiered Electricity Bill]{royal}
\small
A power company charges $10$ cents per kilowatt-hour (kWh) for the first $500$ kWh a household uses
in a month, and $16$ cents per kWh for everything above $500$. The graph shows $E(k)$, the monthly
bill in dollars for $k$ kilowatt-hours, for households using up to $1000$ kWh.

\begin{center}
\begin{tikzpicture}
\begin{axis}[
    width=10.6cm, height=4.4cm,
    xlabel={$k$ (kilowatt-hours used)}, ylabel={$E$ (bill, \$)},
    xmin=0, xmax=1000, ymin=0, ymax=140,
    xtick={0,100,200,300,400,500,600,700,800,900,1000},
    ytick={0,20,40,60,80,100,120}, minor y tick num=1,
    grid=both, grid style={linegray!45}, minor grid style={linegray!30},
    axis lines=left,
    tick label style={font=\tiny}, label style={font=\scriptsize}]
  \addplot[royal, very thick] coordinates {(0,0) (500,50) (1000,130)};
  \addplot[keyred, only marks, mark=*, mark size=1.7pt]
    coordinates {(200,20) (500,50) (750,90) (1000,130)};
\end{axis}
\end{tikzpicture}
\end{center}

\begin{enumerate}[leftmargin=1.6em, itemsep=6pt, topsep=2pt, start=6]

\item Read the graph and fill in both rows. Then give the boundary point of $E$ and say what it means
on a customer's bill.
\smallskip

\renewcommand{\arraystretch}{1.3}
\begin{tabular}{@{}|l|c|c|c|c|@{}}
\hline
$k$ (kWh) & $200$ & $500$ & $750$ & $1000$ \\ \hline
$E(k)$ (\$) & \ans{$20$} & \ans{$50$} & \ans{$90$} & \ans{$130$} \\ \hline
Which piece? & \ans{1st} & \ans{1st} & \ans{2nd} & \ans{2nd} \\
\hline
\end{tabular}

\ansline{Boundary point: $k=500$ kWh. Every kilowatt-hour up to and including the $500$th costs $10$}\\
\ansline{cents; from the $501$st onward each one costs $16$ cents. The cheap rate does not stop applying.}\\

\item Name the interval on which $E$ is increasing. Then say whether $E$ is continuous at $k=500$,
and what the graph does there instead of breaking.

\ansline{$E$ is increasing on all of $[0,1000]$ --- more electricity always costs more.}\\
\ansline{It is continuous at $k=500$: both rules give \$50 there, so there is no gap and no open circle.}\\
\ansline{Instead of breaking, the graph turns a \textbf{corner} --- it gets steeper, because the rate went up.}\\

\item Give the absolute maximum and the absolute minimum of $E$, each with a value and a location.
Then give the zero of $E$ and say what it means about a household.

\ansline{Absolute maximum \$130 at $k=1000$; absolute minimum \$0 at $k=0$. Both sit at endpoints of}\\
\ansline{the domain, because the graph increases the whole way across.}\\
\ansline{Zero at $k=0$ (also the $y$-intercept $(0,0)$): a household that used no electricity owes nothing.}\\

\item A customer writes to the company: \emph{``I used $500$ kWh in April and \$50 was fair. In May I
used $1000$ kWh --- exactly double --- and you charged me \$130. Doubling my usage should have
doubled my bill. You are overcharging me.''} Settle the complaint using the two pieces. Is the
company's arithmetic wrong, or is the customer's expectation wrong?

\ansline{The arithmetic is right; the expectation is wrong. Doubling the input doubles the output only}\\
\ansline{when one single rate covers the whole range --- and here it does not. The first $500$ kWh cost}\\
\ansline{$500(\$0.10)=\$50$, and the \emph{second} $500$ cost $500(\$0.16)=\$80$, for \$130 in total. The extra}\\
\ansline{\$30 is exactly $500$ kWh $\times$ the $6$-cent difference. That is what a tiered rate is designed to do.}\\

\end{enumerate}
\end{scenariobox}

\vspace{0.10in}

\boxguard[24]
\begin{scenariobox}[A First Look Ahead]{goldacc}
\small
\begin{enumerate}[leftmargin=1.6em, itemsep=6pt, topsep=2pt, start=10]

\item Every function family in this unit describes some situations well and others badly. For each
story below, name the family you would use --- \emph{linear}, \emph{quadratic}, \emph{exponential},
or \emph{piecewise} --- and give the one characteristic of the graph that gives it away (a constant
rate of change, a turning point, a horizontal asymptote, a boundary point, and so on). Answer in
plain English; Lesson 2.7 will give you the rest of the vocabulary.
\smallskip

\renewcommand{\arraystretch}{1.4}
\begin{tabularx}{\linewidth}{@{}X p{2.5cm} p{3.6cm}@{}}
\rowcolor{mist}
\textbf{The situation} & \textbf{Family} & \textbf{What gives it away} \\
\hline
A gym charges a \$50 joining fee plus \$30 a month. & \ans{linear} & \ans{one constant rate of
    change} \\
A colony of bacteria doubles every hour. & \ans{exponential} & \ans{a constant multiplier, not a
    constant difference} \\
A ball is thrown straight up; it rises, turns, and falls back. & \ans{quadratic} & \ans{one turning
    point, symmetric} \\
A tax is $10\%$ on the first \$10{,}000 earned and $20\%$ on everything above that. &
    \ans{piecewise} & \ans{a boundary point, then a corner} \\
\end{tabularx}

\smallskip
Then answer: two of these four could be confused with each other if you only saw part of the graph.
Which two, and what would you have to look at to tell them apart?

\ansline{The gym (linear) and the tax (piecewise). Below \$10{,}000 the tax graph is a straight line and}\\
\ansline{looks exactly linear --- you have to look \emph{past the boundary point} to see the slope change.}\\
\ansline{(Quadratic and exponential is also fair: both rise, and only the far ends tell them apart.)}\\

\end{enumerate}
\end{scenariobox}

\vspace{0.10in}

\boxguard[26]
\begin{extensionbox}
\small
\textbf{A graph you already know, hiding in two pieces.} Consider the function
\[
  A(x) = \begin{cases}
    -x & \text{if } x < 0,\\[2pt]
    x & \text{if } x \ge 0.
  \end{cases}
\]

\begin{enumerate}[label=(\alph*), leftmargin=1.6em, itemsep=5pt, topsep=4pt]
  \item Complete the table. Watch $x=0$ carefully --- which of the two pieces owns it?
\smallskip

\renewcommand{\arraystretch}{1.3}
\begin{tabular}{@{}|l|c|c|c|c|c|c|c|@{}}
\hline
$x$ & $-3$ & $-2$ & $-1$ & $0$ & $1$ & $2$ & $3$ \\ \hline
$A(x)$ & \ans{$3$} & \ans{$2$} & \ans{$1$} & \ans{$0$} & \ans{$1$} &
    \ans{$2$} & \ans{$3$} \\
\hline
\end{tabular}

  \item Describe the shape of this graph in one sentence. Is $A$ continuous at $x=0$? What does the
        graph do there?

\ansline{A ``V'' with its point at the origin --- this is the absolute value function, $A(x)=|x|$.}\\
\ansline{Continuous at $x=0$: the second piece owns $x=0$ and both rules give $0$ there, so no gap.}\\
\ansline{The graph turns a sharp \textbf{corner} at the origin --- and a corner is not a break.}\\

  \item Give the interval on which $A$ is decreasing and the interval on which it is increasing, the
        zeros of $A$, the absolute minimum with its location, and the end behavior at both ends
        (Lesson 2.5).

\ansline{Decreasing on $(-\infty,0]$; increasing on $[0,\infty)$. One zero, at $x=0$.}\\
\ansline{Absolute minimum $0$, at $x=0$ --- the corner. There is no absolute maximum.}\\
\ansline{End behavior: the outputs climb without bound at \emph{both} ends, so there is no asymptote.}\\
\end{enumerate}
\end{extensionbox}

\vspace{0.10in}

\begin{remindbox}
\small
\textbf{Next: Lesson 2.7 --- Choosing and Comparing Models in Context} (\texttt{AFDA.AF.1c, e, g};
\texttt{AFDA.AF.2h}). Item 10 asked you to pick a function family from a story. Next lesson turns that
into a method: given data or a situation, decide \emph{which} family models it, defend the choice
using the characteristics you have spent this whole unit learning to read, and say what the model can
and cannot tell you.
\end{remindbox}

\end{document}
